% 06-vmap.tex
\section{Vmap: Batched Transforms and Transform Dispatch}
\label{sec:vmap}

% Coverage: DispatchKey Vmap at offset 9 and priority; PushVmapMode/PopVmapMode level; Vmap->Autocast->Autograd->Backend

Vmap turns a function written for a single sample into a function that runs over a batch without rewriting its body. The public entry point is \path{tensorplay.func.vmap} (\path{tensorplay/\_transforms/apis.py:44}) which delegates to \path{tensorplay.\_transforms.vmap.vmap\_impl} (\path{tensorplay/\_transforms/vmap.py:388}) and ultimately to the native layer in \texttt{p10}: \texttt{push\_vmap} / \texttt{make\_batched} / \texttt{unwrap\_at\_level} (\path{p10/include/TransformDispatch.h:37}, \path{p10/src/TransformDispatch.cpp:24}). Every operator that participates in vmap has a per-backend batch rule registered under the \texttt{Vmap*} dispatch key (\path{p10/src/BatchingKernels.cpp:1223}), and operators without a rule fail fast in generated code (\path{tools/codegen/gen\_api.py:617}). The dispatch key for vmap outranks all other keys (\path{p10/include/DispatchKey.h:41}), so \texttt{TransformDispatch} unwraps before \texttt{Autocast} or \texttt{Autograd} ever see the operands. This section documents the mechanism end to end, with line-precise citations.

\subsection{Dispatch Key at Offset 9 and Priority}
\label{sec:vmapkey}

\path{p10/include/DispatchKey.h:15--49} defines 13 keys with fixed numeric values. The vmap family is placed at offset~9:

\begin{lstlisting}[caption={\path{DispatchKey.h:15} --- 13 keys, Vmap at offset~9.}, label=lst:vmapkeys]
enum class DispatchKey : uint8_t {
  CPU=0, CUDA=1, Vulkan=2,
  AutogradCPU=3, AutogradCUDA=4, AutogradVulkan=5,
  AutocastCPU=6, AutocastCUDA=7, AutocastVulkan=8,
  VmapCPU=9, VmapCUDA=10, VmapVulkan=11,
  Composite=12, EndOfKeys=13
};
constexpr size_t kBackendKeyCount=3, kAutogradKeyOffset=3, kAutocastKeyOffset=6, kVmapKeyOffset=9;
\end{lstlisting}

Helpers \texttt{toVmapKey}, \texttt{is\_vmap\_key}, and \texttt{toBackendKey} are \texttt{constexpr} (\path{DispatchKey.h:64--101}). \texttt{DispatchKeySet} (\path{DispatchKey.h:125--187}) is a \texttt{uint32\_t} bitmask with \path{add/has} and \path{highest\_priority\_key()} which scans from the most-significant set bit downward; because vmap occupies 9--11 it always wins over autocast (6--8), autograd (3--5), and backend (0--2). Table~\ref{tab:vmappriority} makes the ordering explicit.

\begin{table}[H]
\centering
\caption{Dispatch priority. \path{highest\_priority\_key()} is the dispatch decision; vmap is top.}
\label{tab:vmappriority}
\small
\begin{tabular}{clcl}
\toprule
Rank & Numeric & Key family & Defined at \\
\midrule
4 & 9--11 & \path{VmapCPU / VmapCUDA / VmapVulkan} & \path{DispatchKey.h:41}, \texttt{kVmapKeyOffset=9} \\
3 & 6--8  & \path{AutocastCPU / AutocastCUDA / AutocastVulkan} & \path{DispatchKey.h:28}, \path{kAutocastKeyOffset=6} \\
2 & 3--5  & \path{AutogradCPU / AutogradCUDA / AutogradVulkan} & \path{DispatchKey.h:22}, \path{kAutogradKeyOffset=3} \\
1 & 0--2  & \path{CPU / CUDA / Vulkan} (dense backends) & \path{DispatchKey.h:17} \\
0 & 12    & \texttt{Composite} (backend-neutral fallback) & \path{DispatchKey.h:36} \\
\bottomrule
\end{tabular}
\end{table}

\path{DispatchKey.h:81} defines \texttt{is\_vmap\_key} as the half-open interval \path{[kVmapKeyOffset, kVmapKeyOffset+kBackendKeyCount)}; \path{toVmapKey(backend) = backend + 9} and \path{toBackendKey(vmap) = vmap - 9} (\path{DispatchKey.h:64}, \path{DispatchKey.h:94}). \path{DispatchKeySet::remove\_vmap()} (\path{DispatchKey.h:177}) clears exactly those three bits, used when re-dispatching below the vmap layer. The \texttt{DispatchTable} (\path{p10/include/Dispatcher.h:49}) has \path{array<atomic<void*>,13>} slots so each vmap key has its own column; \path{OperatorHandle::getKernel} consults \texttt{Composite} only for backend keys, never for vmap keys, so a missing vmap kernel is a hard miss rather than a silent fallback.

\subsection{TensorImpl Batching State: \texttt{transform\_value\_} Wrapper}
\label{sec:batchstate}

Batching state lives in \texttt{TensorImpl} (\path{p10/include/TensorImpl.h:28--84}) as three fields, separate from storage, version, and sparse metadata:

\begin{lstlisting}[caption={\path{TensorImpl.h:79} --- batching fields.}, label=lst:batchfields]
std::shared_ptr<TensorImpl> transform_value_;
int64_t transform_batch_dim_ = -1;
int64_t transform_level_ = -1;
bool is_batched() const { return transform_value_ != nullptr; }
int64_t batch_dim() const { return transform_batch_dim_; }
int64_t batch_level() const { return transform_level_; }
int64_t batch_size() const {
  return is_batched() ? transform_value_->size(static_cast<size_t>(transform_batch_dim_)) : 0;
}
\end{lstlisting}

The wrapper constructor (\path{p10/include/TensorImpl.h:103}, \path{p10/src/TensorImpl.cpp:65}) takes a physical \path{shared\_ptr<TensorImpl> value} plus logical \path{sizes/strides}:

\begin{lstlisting}[caption={\path{TensorImpl.cpp:65} --- wrapper construction.}, label=lst:wrapperctor]
TensorImpl::TensorImpl(std::shared_ptr<TensorImpl> transform_value,
                       const std::vector<int64_t>& sizes,
                       const std::vector<int64_t>& strides)
  : storage_offset_(0), sizes_and_strides_(sizes, strides),
    dtype_(transform_value ? transform_value->dtype() : DType::Undefined),
    device_(transform_value ? transform_value->device() : Device(DeviceType::CPU)),
    is_contiguous_(sizes_and_strides_.is_contiguous()),
    transform_value_(std::move(transform_value)) {
  shared_state_ = std::make_shared<SharedState>();
}
\end{lstlisting}

The public shape is the logical shape (without the batch dimension); the physical shape lives in \texttt{transform\_value\_}. Views and copies preserve the distinction: \path{TensorImpl(const TensorImpl\& other)} copies \path{transform\_value\_/batch\_dim\_/level\_} (\path{TensorImpl.cpp:87}), while \path{share\_version\_counter} is \emph{not} shared for batched wrappers---batching is orthogonal to versioning. The logical \path{sizes\_and\_strides\_} is recomputed per result by the batch rule; the physical tensor is never aliased as the public tensor.

\path{TensorImpl::key\_set()} (\path{p10/include/TensorImpl.h:135--147}) computes the dispatch key on the fly and gives vmap priority by early return:

\begin{lstlisting}[caption={\path{TensorImpl.h:135} --- computed key set, vmap first.}, label=lst:keysetvmap]
DispatchKeySet key_set() const {
  DispatchKey backend = computeDispatchKey(device_);
  DispatchKeySet ks;
  if (is_batched()) {
    ks.add(toVmapKey(backend));
    return ks;
  }
  ks.add(backend);
  if (!inference_tensor_) {
    ks.add(toAutogradKey(backend));
  }
  return ks;
}
\end{lstlisting}

When \texttt{is\_batched()} the set contains only the \texttt{Vmap*} bit, so \path{dispatchKeyForTensorArgs} (\path{p10/include/Tensor.h:417}) will select a vmap key if any argument is batched, regardless of how many unbatched arguments are present. Table~\ref{tab:batchfields} summarizes the field semantics.

\begin{table}[H]
\centering
\caption{\texttt{TensorImpl} batching members. Logical vs.\ physical split.}
\label{tab:batchfields}
\small
\begin{tabular}{lll}
\toprule
Field & Type & Meaning \\
\midrule
\texttt{transform\_value\_} & \path{shared\_ptr<TensorImpl>} & Physical tensor (with batch dim); \texttt{nullptr} means not batched \\
\path{transform\_batch\_dim\_} & \texttt{int64\_t} & Position of batch dim inside the physical shape; $-1$ when unbatched \\
\texttt{transform\_level\_} & \texttt{int64\_t} & Which vmap nesting level owns this wrapper; distinguishes nested vmaps \\
\path{sizes\_and\_strides\_} & \texttt{SizesAndStrides} & Logical shape (public, without batch dim) \\
\bottomrule
\end{tabular}
\end{table}

The thin handle \texttt{Tensor} (\path{p10/include/Tensor.h:167}) forwards \path{is\_batched()/batch\_dim()/batch\_level()/transform\_value()} to the impl with a null check, so callers never touch \texttt{TensorImpl} directly.

\subsection{Transform Stack: Push/Pop and Levels}
\label{sec:transformstack}

Nested vmaps are disambiguated by a per-thread level stack (\path{p10/src/TransformDispatch.cpp:11}, \path{p10/include/TransformDispatch.h:27}):

\begin{lstlisting}[caption={\path{TransformDispatch.h:27} --- layer stack.}, label=lst:layerstack]
enum class Randomness : uint8_t { Error=0, Same=1, Different=2 };
enum class Kind : uint8_t { Vmap=0, Grad=1, Jvp=2, Functionalize=3 };
struct Layer {
  Kind kind = Kind::Vmap;
  int64_t level = -1;
  int64_t batch_size = 0;
  Randomness randomness = Randomness::Error;
  bool previous_grad_mode = true;
  bool previous_forward_grad_mode = true;
  bool add_back_views = false;
};
P10_API int64_t push_vmap(int64_t batch_size, Randomness randomness);
P10_API Layer pop_layer();
P10_API std::optional<Layer> current_layer();
P10_API std::vector<Layer> layer_stack();
\end{lstlisting}

The implementation (\path{p10/src/TransformDispatch.cpp:11--57}) is a thread-local vector plus a monotonic counter:

\begin{lstlisting}[caption={\path{TransformDispatch.cpp:11} --- stack and levels.}, label=lst:pushpop]
namespace { thread_local std::vector<Layer> layers; thread_local int64_t next_level = 0; }

int64_t push_vmap(int64_t batch_size, Randomness randomness) {
  if (batch_size < 0) TP_THROW(ValueError, "vmap batch size must be non-negative");
  Layer layer; layer.kind = Kind::Vmap; layer.level = next_level++;
  layer.batch_size = batch_size; layer.randomness = randomness;
  layers.push_back(layer); return layer.level;
}
Layer pop_layer() {
  if (layers.empty()) TP_THROW(RuntimeError, "cannot pop an empty transform stack");
  Layer layer = layers.back(); layers.pop_back(); return layer;
}
\end{lstlisting}

Levels are never reused within a thread: even after \texttt{pop\_layer()}, \texttt{next\_level} keeps incrementing. This prevents ABA confusion when a batched tensor outlives its layer (e.g., captured in a closure). \path{is\_batched\_at\_level(value, level)} (\path{TransformDispatch.cpp:109}) checks both the pointer and the exact level, so an outer vmap's tensor is invisible to an inner vmap's rules.

The Python guard \path{vmap\_increment\_nesting} (\path{tensorplay/\_transforms/vmap.py:142}) pushes, yields the level, and pops in a \texttt{finally} block that also restores a \texttt{ContextVar} depth counter (\path{vmap.py:23}, \path{vmap.py:145}). The names \texttt{PushVmapMode}/\texttt{PopVmapMode} in the task description map to \texttt{push\_vmap}/\texttt{pop\_layer} in code (\path{TransformDispatch.h:37}, \path{TransformDispatch.cpp:24}); the ``Mode'' wording reflects the RAII guard \path{vmap\_increment\_nesting} that pushes and pops a mode, while the native symbols are plain functions. The C++ binding (\path{src/bindings/python/Transforms.cpp:18}) exposes \path{\_transform\_push\_vmap} / \texttt{\_transform\_pop} / \path{\_transform\_current} directly, with string-to-enum parsing for randomness (\path{Transforms.cpp:7}). No Python-level mode flag is needed beyond the C++ stack; the stack itself is the ground truth.

Operational invariants:
\begin{enumerate}[leftmargin=1.2em]
\item \texttt{push\_vmap} returns a fresh level; \path{make\_batched(value, dim, level)} tags the wrapper with that level (\path{TransformDispatch.cpp:66}).
\item \path{unwrap\_at\_level(value, level)} returns \texttt{(physical, bdim)} only when \path{value.batch\_level()==level} (\path{TransformDispatch.cpp:93}); otherwise it returns \texttt{(value, nullopt)} unchanged.
\item \path{dispatch\_key\_for\_random(backend)} returns \texttt{toVmapKey(backend)} exactly when a vmap layer is on top of the stack (\path{TransformDispatch.cpp:59}), so random factories inside vmap re-dispatch through the vmap key even though they have no tensor argument to carry it.
\end{enumerate}

\subsection{TransformDispatch Unwrapping Logic}
\label{sec:unwrap}

\texttt{TransformDispatch} provides the primitives that batch rules use to inspect and rebuild batched tensors (\path{TransformDispatch.h:44}, \path{TransformDispatch.cpp:66--140}):

\begin{lstlisting}[caption={\path{TransformDispatch.cpp:66} --- make and unwrap.}, label=lst:makeunwrap]
Tensor make_batched(const Tensor& value, int64_t dim, int64_t level) {
  if (!value.defined()) TP_THROW(ValueError, "cannot batch an undefined tensor");
  const int64_t bdim = normalize_dim(dim, value.dim());
  const auto sizes = static_cast<std::vector<int64_t>>(value.shape());
  const auto strides = value.strides();
  std::vector<int64_t> logical_sizes, logical_strides;
  for (int64_t d = 0; d < (int64_t)sizes.size(); ++d) if (d != bdim) {
    logical_sizes.push_back(sizes[d]); logical_strides.push_back(strides[d]);
  }
  auto impl = std::make_shared<TensorImpl>(value.unsafeGetTensorImpl(),
                                           logical_sizes, logical_strides);
  impl->set_transform_value(value.unsafeGetTensorImpl(), bdim, level);
  return Tensor(std::move(impl));
}
std::tuple<Tensor, std::optional<int64_t>> unwrap_at_level(const Tensor& value, int64_t level) {
  if (!value.defined() || !value.is_batched() || value.batch_level() != level)
    return {value, std::nullopt};
  return {value.transform_value(), value.batch_dim()};
}
Tensor move_batch_dim(const Tensor& value, int64_t dim) {
  if (!value.is_batched()) return value;
  const int64_t from = normalize_dim(value.batch_dim(), value.dim()+1);
  const int64_t to = normalize_dim(dim, value.dim()+1);
  if (from == to) return value;
  std::vector<int64_t> perm; // build permutation that moves from->to
  // ... permute physical tensor
  return value.transform_value().permute(perm);
}
\end{lstlisting}

\texttt{actual\_dim} (\path{TransformDispatch.cpp:113}) maps a public dimension (logical) to a physical dimension by skipping the batch position: \path{actual = public $<$ bdim ? public : public+1}. Every batch rule that takes a dim argument (permute, transpose, squeeze, sum, etc.) applies this translation before calling the underlying kernel, and then adjusts the result's \texttt{bdim} to reflect dims inserted or removed before the batch position.

The unwrapping contract is uniform across \texttt{BatchingKernels.cpp}: each rule calls \texttt{unwrap\_operand} (\path{BatchingKernels.cpp:34}) which asserts the operand's level matches an active layer via \texttt{layer\_for(level)} (\path{BatchingKernels.cpp:22}), then dispatches to \texttt{unwrap\_at\_level}. If no operand is batched at the current level, the rule throws \texttt{RuntimeError}---the dispatcher should not have routed there. This fail-fast behavior catches codegen mismatches immediately.

\begin{figure}[H]
\centering
\begin{tikzpicture}[every node/.style={font=\scriptsize}, node distance=0.55cm]
\node[libbox, fill=blue!8, draw=tpblue, minimum width=2.8cm, minimum height=0.9cm] (py) {\texttt{vmap(f)(x)}};
\node[libbox, fill=orange!8, draw=tporange, below=of py, minimum width=2.8cm, minimum height=0.9cm] (push) {\texttt{push\_vmap(B, rand)}};
\node[libbox, fill=purple!8, draw=tppurple, below=of push, minimum width=2.8cm, minimum height=0.9cm] (make) {\path{make\_batched(x, 0, lvl)}};
\node[libbox, fill=green!8, draw=tpgreen, below=of make, minimum width=2.8cm, minimum height=0.9cm] (dispatch) {\texttt{Tensor::add} $\rightarrow$ \texttt{VmapCPU}};
\node[libbox, fill=gray!8, draw=gray, below=of dispatch, minimum width=2.8cm, minimum height=0.9cm] (rule) {\texttt{batch\_add}: unwrap, \texttt{call\_next}};
\node[libbox, fill=blue!8, draw=tpblue, below=of rule, minimum width=2.8cm, minimum height=0.9cm] (wrap) {\path{make\_batched(result, 0, lvl)}};
\node[libbox, fill=orange!8, draw=tporange, below=of wrap, minimum width=2.8cm, minimum height=0.9cm] (pop) {\texttt{pop\_layer()} + \texttt{movedim}};
\draw[arrow] (py) -- (push); \draw[arrow] (push) -- (make); \draw[arrow] (make) -- (dispatch);
\draw[arrow] (dispatch) -- (rule); \draw[arrow] (rule) -- (wrap); \draw[arrow] (wrap) -- (pop);
\node[right=0.6cm of dispatch, align=left, font=\tiny, text=gray] (note) {\texttt{key\_set()} returns\\\texttt{VmapCPU} only};
\draw[dashed, color=gray] (note) -- (dispatch);
\end{tikzpicture}
\caption{Vmap execution: push level, wrap inputs, dispatch through Vmap key, unwrap inside batch rule, rewrap result, pop and place output dim.}
\label{fig:vmapflow}
\end{figure}

\texttt{unwrap\_all} (\path{TransformDispatch.cpp:101}) strips all batch wrappers by looping; it is used by \texttt{debug\_unwrap} and by composites that need the physical tensor for metadata queries. \texttt{move\_batch\_dim} (\path{TransformDispatch.cpp:126}) permutes the physical tensor so the batch dim lands at the requested public position, used when \texttt{out\_dims != 0}.

\subsection{BatchingKernels: Per-Op Hand-Written Rules vs.\ Loop Fallback}
\label{sec:batchingkernels}

\path{p10/src/BatchingKernels.cpp} is 1229 lines of batch rules, each registered under \texttt{VmapCPU} and \texttt{VmapCUDA} (\path{BatchingKernels.cpp:1223}):

\begin{lstlisting}[caption={\path{BatchingKernels.cpp:1223} --- dual registration.}, label=lst:batchreg]
TENSORPLAY_LIBRARY_IMPL(VmapCPU, NativeBatchRulesCPU) { register_batch_rules(m); }
TENSORPLAY_LIBRARY_IMPL(VmapCUDA, NativeBatchRulesCUDA) { register_batch_rules(m); }
\end{lstlisting}

\path{register\_batch\_rules} (\path{BatchingKernels.cpp:1135}) binds ~70 ops. The file has no generic per-slice loop fallback; every supported op has a hand-written rule that preserves vectorization. Unsupported ops are not silently looped---the generated \texttt{Tensor} method (\path{tools/codegen/gen\_api.py:617}) checks \path{getKernel(vmap\_key)} and throws \path{No native batch rule registered} (\path{gen\_api.py:617}--\path{gen\_api.py:624}) if absent, which is the correct fail-fast for a teaching framework where implicit loops would hide performance cliffs.

\subsubsection{Rule templates}

Three templates cover most ops (\path{BatchingKernels.cpp:306--325}):

\begin{lstlisting}[caption={\path{BatchingKernels.cpp:306} --- unary / binary templates.}, label=lst:templates]
template <typename Function> Tensor unary_impl(const Tensor& input, Function&& call) {
  Operand operand = unwrap_operand(input);
  if (!operand.bdim.has_value()) TP_THROW(RuntimeError, "unary batch rule received an unbatched operand");
  Tensor result = call(operand.value);
  return make_batched(result, *operand.bdim, operand.level);
}
template <typename Function> Tensor binary_impl(const Tensor& left, const Tensor& right, Function&& call) {
  Operand left_operand = unwrap_operand(left); Operand right_operand = unwrap_operand(right);
  int64_t level = -1; auto values = aligned_values(left_operand, right_operand, level);
  Tensor result = call(values.first, values.second);
  return make_batched(result, 0, level);
}
Tensor unary(const char* op, const Tensor& input) {
  return unary_impl(input, [&](const Tensor& v){ return call_next<Tensor,const Tensor&>(op, v, v); });
}
Tensor binary(const char* op, const Tensor& left, const Tensor& right) {
  return binary_impl(left, right, [&](const Tensor& a, const Tensor& b){
    return call_next<Tensor,const Tensor&,const Tensor&>(op, a, a, b); });
}
\end{lstlisting}

\texttt{call\_next} (\path{BatchingKernels.cpp:119}) re-dispatches below the vmap layer via \texttt{DispatchStub::call} with the backend key derived from the unwrapped tensors, so the underlying CPU/CUDA kernel runs on contiguous physical memory. \texttt{call\_device} (\path{BatchingKernels.cpp:131}) is the factory variant that dispatches by explicit \texttt{Device}.

Key helpers:
\begin{itemize}[leftmargin=1.2em]
\item \path{move\_to\_front(value, dim)} (\path{BatchingKernels.cpp:56}) permutes batch dim to 0, enabling a single batched kernel to handle any \texttt{bdim}.
\item \path{expand\_unbatched(value, batch\_size, target\_shape)} (\path{BatchingKernels.cpp:68}) broadcasts an unbatched operand by \path{unsqueeze(0).expand([B, ...])} so binary ops need not branch on mixed batched/unbatched in the compute kernel.
\item \path{aligned\_values(left, right, level)} (\path{BatchingKernels.cpp:88}) moves batched inputs to front and expands unbatched ones, returning two tensors whose dim~0 is the batch. The result is always wrapped at dim~0.
\item \path{logical\_shape(operand)} (\path{BatchingKernels.cpp:77}) erases the batch dim from the physical shape, used to size expansions.
\end{itemize}

\subsubsection{Operator families}

Table~\ref{tab:batchfamilies} groups the registered rules.

\begin{table}[H]
\centering
\caption{Batch rule families in \path{BatchingKernels.cpp:1135}. Each row is a \texttt{library.impl} entry.}
\label{tab:batchfamilies}
\small
\begin{tabular}{lll}
\toprule
Family & Examples & Strategy \\
\midrule
Pointwise unary & \path{neg, abs, exp, log, sin, cos, sqrt, sigmoid, relu} & \texttt{unary\_impl} preserves \texttt{bdim} (\path{BatchingKernels.cpp:936}) \\
Pointwise binary & \path{mul.Tensor, div.Tensor, maximum, minimum} & \texttt{binary\_impl} via \texttt{aligned\_values}, result at 0 \\
Binary with alpha & \path{add.Tensor, sub.Tensor} & \texttt{binary\_alpha}, same alignment \\
Scalar & \path{add.Scalar, mul.Scalar, pow.Tensor\_Scalar} & \texttt{scalar}/\texttt{scalar\_alpha}, preserves \texttt{bdim} \\
Reductions & \path{sum, sum.dim\_IntList} & Rewrites dims past \texttt{bdim}, adjusts result \texttt{bdim} (\path{BatchingKernels.cpp:381}) \\
Shape & \path{view, permute, transpose, reshape, expand, squeeze, unsqueeze} & Translates public dims to physical, adjusts result \texttt{bdim} \\
Gather/scatter & \path{select, slice, narrow, index\_select, cat, stack} & Shifts gather dim past \texttt{bdim}; \path{cat/stack} align then \texttt{call\_next} at \texttt{dim+1} \\
Linear algebra & \path{mm, matmul, bmm, linear} & Align to front, call \path{bmm/matmul}; \texttt{linear} transposes weight (\path{BatchingKernels.cpp:834}) \\
Random & \path{rand, randn, randint, randperm, rand\_like} & \texttt{random\_factory} branches on \texttt{Randomness} (see \S\ref{sec:randomness}) \\
\bottomrule
\end{tabular}
\end{table}

Concrete dim adjustment appears in every shape rule. For example, \texttt{sum\_dim} (\path{BatchingKernels.cpp:400}):

\begin{lstlisting}[caption={\path{BatchingKernels.cpp:400} --- sum adjusts batch dim.}, label=lst:sumdim]
Tensor sum_dim(const Tensor& input, const std::vector<int64_t>& dims, bool keepdim, DType dtype) {
  Operand operand = unwrap_operand(input);
  const int64_t old_bdim = *operand.bdim; const int64_t level = operand.level;
  std::vector<int64_t> actual_dims;
  for (int64_t dim : dims) {
    const int64_t public_dim = normalize_dim(dim, input.dim());
    const int64_t actual = public_dim < old_bdim ? public_dim : public_dim + 1;
    actual_dims.push_back(actual);
  }
  Tensor result = call_next<Tensor,const Tensor&,const std::vector<int64_t>&,bool,DType>(
    "sum.dim_IntList", operand.value, operand.value, actual_dims, keepdim, dtype);
  int64_t result_bdim = old_bdim;
  if (!keepdim) result_bdim -= count_if(actual_dims.begin(), actual_dims.end(),
                                        [&](int64_t d){ return d < old_bdim; });
  return make_batched(result, result_bdim, level);
}
\end{lstlisting}

\texttt{permute} (\path{BatchingKernels.cpp:445}) builds a physical permutation that interleaves the batch dim at its original position, then calls \texttt{permute} on the physical tensor and rewraps at \texttt{new\_bdim}. \texttt{cat} (\path{BatchingKernels.cpp:767}) aligns a list via \texttt{align\_tensor\_list} (\path{BatchingKernels.cpp:736}) which moves every batched input's \texttt{bdim} to 0 and expands unbatched ones, then calls \texttt{cat} at \texttt{public\_dim+1} and wraps at 0. \texttt{mm} (\path{BatchingKernels.cpp:793}) aligns then calls \texttt{bmm} so a batch of matrix multiplies becomes a single batched kernel rather than a loop.

No rule falls back to a per-slice loop. This is intentional: a loop would be correct but would defeat the purpose of a batched kernel (one launch vs.\ $B$ launches) and would hide missing coverage. When an op lacks a rule, the generated wrapper throws immediately, and the fix is to add a rule to \texttt{BatchingKernels.cpp} following the \texttt{unary}/\texttt{binary} pattern.

\subsection{Python Surface: \path{tensorplay.func.vmap}}
\label{sec:pythonsurface}

The Python API is a thin validator over the native stack. \path{tensorplay/\_transforms/apis.py:44} is the public \texttt{vmap}:

\begin{lstlisting}[caption={\path{apis.py:44} --- public \texttt{vmap} surface.}, label=lst:pyvmapapi, style=pystyle]
def vmap(func, in_dims=0, out_dims=0, randomness="error", *, chunk_size=None):
    if not callable(func): raise TypeError(f"vmap expected a callable, got {type(func)!r}")
    _check_randomness_arg(randomness)
    @functools.wraps(func)
    def wrapped(*args, **kwargs):
        return vmap_impl(func, in_dims, out_dims, randomness, chunk_size, *args, **kwargs)
    return wrapped
\end{lstlisting}

\path{tensorplay/\_transforms/vmap.py:388} validates, flattens the pytree of args, pushes a native layer, wraps inputs, calls the user function, and unwraps outputs:

\begin{lstlisting}[caption={\path{vmap.py:388} --- \texttt{vmap\_impl} and nesting.}, label=lst:vmapimpl, style=pystyle]
def vmap_impl(func, in_dims, out_dims, randomness, chunk_size, *args, **kwargs):
    _check_randomness_arg(randomness)
    batch_size, flat_in_dims, flat_args, args_spec = _process_batched_inputs(in_dims, args, func)
    if chunk_size is not None:
        return _chunked_vmap(func, flat_in_dims,
            _get_chunked_inputs(flat_args, flat_in_dims, batch_size, chunk_size),
            args_spec, out_dims, randomness, **kwargs)
    return _flat_vmap(func, batch_size, flat_in_dims, flat_args, args_spec, out_dims, randomness, **kwargs)

def _flat_vmap(func, batch_size, flat_in_dims, flat_args, args_spec, out_dims, randomness, **kwargs):
    with vmap_increment_nesting(batch_size, randomness) as level:
        batched_args = _create_batched_inputs(flat_in_dims, flat_args, level, args_spec)
        outputs = func(*batched_args, **kwargs)
        return _unwrap_batched(outputs, out_dims, level, batch_size, func)

@contextlib.contextmanager
def vmap_increment_nesting(batch_size, randomness):
    level = tensorplay._C._transform_push_vmap(batch_size, randomness)
    token = _vmap_depth.set(_vmap_depth.get() + 1)
    try: yield level
    finally:
        try: tensorplay._C._transform_pop()
        finally: _vmap_depth.reset(token)
\end{lstlisting}

\path{\_process\_batched\_inputs} (\path{vmap.py:68}) flattens both \texttt{args} and \texttt{in\_dims} via \texttt{tree\_flatten} / \path{\_broadcast\_to\_and\_flatten}, normalizes negative dims, checks that every \texttt{in\_dim != None} corresponds to a \texttt{Tensor}, and verifies all mapped tensors share the same batch size (\path{vmap.py:50}). \path{\_create\_batched\_inputs} (\path{vmap.py:124}) calls \path{tensorplay.\_C.\_transform\_make\_batched(arg, dim, level)} for each mapped leaf, leaving unmapped leaves untouched.

Output handling (\path{vmap.py:190}) follows the input side: it flattens outputs, checks \texttt{out\_dims} compatibility, unwraps each tensor leaf with \path{\_transform\_unwrap(output, level)}, and moves the physical batch dim to the requested public position via \path{tensorplay.movedim(unwrapped, batch\_dim, target\_dim)} (\path{vmap.py:242}). Non-tensor leaves must have \texttt{out\_dim=None}; otherwise a \texttt{ValueError} is raised.

Three chunking utilities bound peak memory for large batches:
\begin{itemize}[leftmargin=1.2em]
\item \texttt{chunk\_vmap} (\path{apis.py:98}) and \path{vmap(..., chunk\_size=N)} (\path{vmap.py:347}) split the batch into pieces via \path{get\_chunked\_inputs} (\path{vmap.py:283}) which uses \texttt{tensor\_split} along the mapped dim, run independent native layers per chunk, and rejoin with \texttt{cat} along \texttt{out\_dim} (\path{vmap.py:309}).
\item \texttt{restore\_vmap} (\path{vmap.py:449}) and \path{wrap\_batched / unwrap\_batched} (\path{vmap.py:478}) expose the same primitives for testing and for \texttt{jvp} which needs to re-enter a sized layer.
\end{itemize}

The C++ binding (\path{src/bindings/python/Transforms.cpp:17}) is six functions: \path{\_transform\_push\_vmap}, \texttt{\_transform\_pop}, \path{\_transform\_current}, \path{\_transform\_make\_batched}, \texttt{\_transform\_unwrap}, and \path{\_transform\_unwrap\_all}. Each is a direct forward to \path{tensorplay::transform::*} with pybind11 type conversions and no extra bookkeeping.

\subsection{Randomness Handling}
\label{sec:randomness}

Random ops inside vmap are ambiguous: should each sample draw independently or should the batch share one draw? TensorPlay requires the caller to choose explicitly via \texttt{randomness} (\path{TransformDispatch.h:14}, \path{vmap.py:507}):

\begin{table}[H]
\centering
\caption{Randomness modes. Every random factory consults the active layer's \texttt{randomness}.}
\label{tab:randomness}
\small
\begin{tabular}{lll}
\toprule
Mode & String & Factory behavior in \texttt{BatchingKernels.cpp} \\
\midrule
\texttt{Error} & \texttt{"error"} & \texttt{check\_randomness} throws (\path{BatchingKernels.cpp:145}) \\
\texttt{Same} & \texttt{"same"} & Draw once, \path{unsqueeze(0).expand([B,...])} and wrap (\path{BatchingKernels.cpp:170}) \\
\texttt{Different} & \texttt{"different"} & Draw with prepended batch dim \texttt{[B, ...shape]} and wrap at 0 \\
\bottomrule
\end{tabular}
\end{table}

Factories (\path{BatchingKernels.cpp:177--275}) implement this branching. For shape-based factories:

\begin{lstlisting}[caption={\path{BatchingKernels.cpp:177} --- random factory.}, label=lst:randfactory]
Tensor random_factory(const char* op, const std::vector<int64_t>& shape,
                      std::optional<DType> dtype, std::optional<Device> device) {
  const Layer layer = current_vmap_layer(); check_randomness(layer);
  const Device target = resolve_random_device(device);
  if (layer.randomness == Randomness::Different) {
    Tensor result = call_device<Tensor,const std::vector<int64_t>&,
                                std::optional<DType>,std::optional<Device>>(
      op, target, prepend_batch_size(shape, layer.batch_size), dtype, device);
    return make_batched(result, 0, layer.level);
  }
  Tensor sample = call_device<Tensor,const std::vector<int64_t>&,
                              std::optional<DType>,std::optional<Device>>(
    op, target, shape, dtype, device);
  return expand_same_random(std::move(sample), layer);
}
\end{lstlisting}

\path{prepend\_batch\_size} (\path{BatchingKernels.cpp:161}) inserts $B$ at dim~0; \path{expand\_same\_random} (\path{BatchingKernels.cpp:170}) does \path{sample.unsqueeze(0).expand([B, ...shape])} so every sample shares the same random values but downstream ops still see a proper batch. \texttt{randperm} (\path{BatchingKernels.cpp:217}) is special: with \texttt{Different} it loops $B$ times and stacks, because a single \texttt{randperm([B,n])} would not yield $B$ independent permutations.

For \texttt{rand\_like} variants (\path{BatchingKernels.cpp:243--304}), the input may itself be batched at the active level. With \texttt{Different}, the factory calls the underlying \texttt{rand\_like} on the physical value and rewraps at the same \texttt{bdim}; with \texttt{Same}, it selects one slice (\texttt{select.int} at \texttt{bdim}, index 0), draws from that slice, and expands back along \texttt{bdim}. Unbatched inputs under \texttt{Same} share one draw expanded to \texttt{[B,...]}; under \texttt{Different} they create a template \texttt{unsqueeze(0).expand} physical tensor and draw there.

The generated dispatch key for random ops (\path{tools/codegen/gen\_api.py:225}) uses \path{dispatch\_key\_for\_random(computeDispatchKey(device))} rather than \path{dispatchKeyForTensorArgs}, so the vmap key is injected even when no tensor argument carries it. At the Python layer, \path{vmap(func, randomness="different")(randn([3]))} therefore routes \texttt{randn} through \texttt{VmapCPU} without any batched input. With the default \texttt{"error"}, any random op inside vmap throws immediately with a message naming the constraint, so silent wrong-batch randomness is impossible.

\subsection{Interaction with Autograd and Redispatch Order}
\label{sec:autogradinteract}

Vmap's numeric priority is load-bearing for correctness: vmap must see operands before autograd records history. \path{TensorImpl::key\_set()} returning only \texttt{Vmap*} when batched, and \path{highest\_priority\_key()} preferring higher numeric keys, together guarantee that \texttt{vmap(grad(f))} and \texttt{grad(vmap(f))} are distinct and both correct.

The generated \texttt{Tensor} methods (\path{tools/codegen/gen\_api.py:609--638}) encode the order Vmap $\rightarrow$ Autocast $\rightarrow$ Autograd $\rightarrow$ Backend:

\begin{lstlisting}[caption={\path{gen\_api.py:609} --- redispatch order in every Tensor method.}, label=lst:redispatchorder]
{
  DispatchKey __tp_key = dispatchKeyForTensorArgs(*this, other);
  if (is_vmap_key(__tp_key)) {
    static const OperatorHandle __tp_handle = Dispatcher::singleton().findHandle("add");
    if (!__tp_handle || !__tp_handle.getKernel(__tp_key))
      TP_THROW(NotImplementedError, "No native batch rule registered for add");
    return DispatchStub<...>::call(__tp_handle, __tp_key, *this, other);
  }
}
{
  static const OperatorHandle __ac_handle = Dispatcher::singleton().findHandle("add");
  DispatchKey __ac_key = toAutocastKey(computeDispatchKey(device()));
  if (autocast::is_enabled(__ac_key) && __ac_handle && __ac_handle.getKernel(__ac_key))
    return DispatchStub<...>::call(__ac_handle, __ac_key, *this, other);
}
if (GradMode::is_enabled()) {
  static const OperatorHandle ag_handle = Dispatcher::singleton().findHandle("add");
  DispatchKey ag_key = toAutogradKey(computeDispatchKey(device()));
  if (ag_handle && ag_handle.getKernel(ag_key))
    return DispatchStub<...>::call(ag_handle, ag_key, *this, other);
}
return detail::redispatch_add(*this, other);
\end{lstlisting}

Table~\ref{tab:vmaporder} summarizes the funnel. The final \path{detail::redispatch\_*} (\path{tools/codegen/gen\_api.py:304}) resolves a backend kernel via \path{toBackendKey(dispatchKeyForTensorArgs(...))} after stripping transform and autograd bits, so the physical compute always lands on \path{CPU/CUDA/Vulkan}.

\begin{table}[H]
\centering
\caption{Redispatch funnel in generated \texttt{Tensor} methods. Each stage either handles the call or falls through.}
\label{tab:vmaporder}
\small
\begin{tabular}{cll}
\toprule
Stage & Key family & Action if kernel present \\
\midrule
1 & \texttt{Vmap*} (\texttt{is\_vmap\_key}) & Unwrap, align, call underlying op, rewrap (\path{BatchingKernels.cpp:306}) \\
2 & \texttt{Autocast*} (\texttt{is\_enabled}) & Cast inputs, call next layer (generated autocast wrapper) \\
3 & \texttt{Autograd*} (\path{GradMode::is\_enabled}) & Record \texttt{Node} with \texttt{SavedVariable} guards, call backend \\
4 & \path{CPU/CUDA/Vulkan} & Compute \\
\bottomrule
\end{tabular}
\end{table}

Because vmap runs first, \texttt{add} inside \texttt{vmap(f)} is dispatched to \texttt{batch\_add} which unwraps to physical tensors, calls \texttt{add} on them (re-entering dispatch at the backend key), and rewraps the result. Autograd therefore sees $B$ batched values as a single batched result and records a batched \texttt{Node}; the backward of a batched \texttt{Node} is itself batched, so \path{Engine::evaluate\_function} needs no vmap-specific changes. Conversely, \texttt{grad(vmap(f))} runs the grad transform outside vmap: the forward records a graph whose intermediate tensors are batched, and the backward replays with the same batch rules.

Composites (\path{p10/src/backend/composite/CompositeCommon.h:21}) reinforce the ordering:

\begin{lstlisting}[caption={\path{CompositeCommon.h:21} --- composites reject active transforms.}, label=lst:composites]
inline bool under_active_transform(const Tensor& t){ return t.unsafeGetTensorImpl() && t.unsafeGetTensorImpl()->is_batched(); }
inline void reject_active_transform(const Tensor& t, const char* op){
  if (under_active_transform(t)) TP_THROW(NotImplementedError, std::string(op),
    " is not supported for tensors inside an active transform (vmap/grad) layer");
}
template <typename Return, typename... Args>
Return redispatch_below_transform(const char* op, Args... args){
  DispatchKey key = dispatchKeyForTensorArgs(args...);
  TP_CHECK(is_vmap_key(key), op, ": expected an active transform key");
  return DispatchStub<Return, Args...>::call(std::string(op), key, std::forward<Args>(args)...);
}
\end{lstlisting}

Plain composites call \path{reject\_active\_transform} because running a backend kernel on the wrapper would silently corrupt batching. Decomposition-style composites call \path{redispatch\_below\_transform} to intentionally route through the vmap layer so inner ops hit their batch rules. Both paths keep the invariant that any op touching a batched tensor must either be a batch rule itself or explicitly opt into re-dispatching through one.

Nested vmaps compose by stacking levels: the inner \texttt{push\_vmap} allocates a new level, its inputs are double-wrapped (outer wrapper around inner wrapper), and each rule's \texttt{unwrap\_at\_level} only peels the level it owns. The outer batch dim remains untouched while the inner rule executes, then the inner result is rewrapped before returning to the outer layer. No special casing for nesting is needed beyond the level check.

\subsection{End-to-End Example and Cost Notes}
\label{sec:vmapcost}

Consider \texttt{f(x)= (x*2).sum()} mapped over 4 samples of shape \texttt{[3]}:

\begin{lstlisting}[caption={Vmap example: one function, batched execution.}, label=lst:vmapexample, style=pystyle]
import tensorplay as tp
def f(x): return (x * 2).sum()
batched = tp.func.vmap(f)(tp.randn([4, 3]))   # in_dims=0, out_dims=0
assert batched.shape == tp.Size([4])
# equivalent to: tp.stack([f(tp.randn([3])) for _ in range(4)]) but one kernel launch per op
\end{lstlisting}

What runs natively:
\begin{enumerate}[leftmargin=1.2em]
\item \texttt{vmap\_impl} validates batch size 4, pushes level $L$ with $B{=}4$, wraps the \texttt{[4,3]} tensor into a batched \texttt{[3]} with \texttt{bdim=0, level=L}.
\item \texttt{x*2} is \texttt{mul.Scalar} with one batched operand: \texttt{batch\_mul\_scalar} unwraps to physical \texttt{[4,3]}, calls \texttt{mul.Scalar} on it (one kernel), rewraps at 0 $\to$ batched \texttt{[3]} at dim~0.
\item \texttt{sum} is \texttt{sum.dim\_IntList}: \texttt{batch\_sum} unwraps, calls \texttt{sum} over the non-batch dims on the physical \texttt{[4,3]}, gets \texttt{[4]}, rewraps at 0 $\to$ batched \texttt{[]} (scalar per sample) with logical shape \texttt{[4]} after output \texttt{movedim}.
\item Pop $L$, move batch dim to \texttt{out\_dims=0} (already there), return \texttt{[4]}.
\end{enumerate}

Cost: each op pays one \texttt{unwrap\_at\_level} + one \texttt{make\_batched} + one underlying kernel. No per-sample Python loop, no extra allocations beyond the wrapper \texttt{TensorImpl} (a few dozen bytes). For a batch of $B$ samples, $N$ ops launch $N$ kernels rather than $N{\times}B$. Chunked vmap (\texttt{chunk\_size}) trades this back for bounded memory by splitting $B$ into pieces and concatenating results.

\subsection{Verification}
\label{sec:vmapverify}

All claims reproduce with local \texttt{grep} and a short Python run:

\begin{lstlisting}[caption={Local checks for vmap.}, label=lst:verifyvmap, style=pystyle]
# Keys and offsets
grep -n "VmapCPU\|kVmapKeyOffset\|is_vmap_key\|highest_priority_key" p10/include/DispatchKey.h
# Batch state and computed key set
grep -n "transform_value_\|is_batched\|batch_dim\|key_set" p10/include/TensorImpl.h | head -n 30
# Stack and primitives
grep -n "push_vmap\|pop_layer\|make_batched\|unwrap_at_level\|move_batch_dim" p10/include/TransformDispatch.h p10/src/TransformDispatch.cpp
# Batch rules and dual registration
grep -n "register_batch_rules\|TENSORPLAY_LIBRARY_IMPL.*Vmap" p10/src/BatchingKernels.cpp | head
grep -n "call_next\|aligned_values\|move_to_front" p10/src/BatchingKernels.cpp | head
# Python surface and codegen order
grep -n "def vmap\|def _flat_vmap\|vmap_increment_nesting\|_check_randomness_arg" tensorplay/_transforms/vmap.py | head
grep -n "is_vmap_key" tools/codegen/gen_api.py | head
# C++ Python bindings
grep -n "_transform_push_vmap\|_transform_make_batched\|_transform_unwrap" src/bindings/python/Transforms.cpp

python3 - <<'PY'
import tensorplay as tp
# basic vmap
def f(x): return (x * 2).sum()
print(tp.func.vmap(f)(tp.randn([4,3])).shape)  # [4]

# out_dims and in_dims
def g(x, y): return x + y
print(tp.func.vmap(g, in_dims=(0, None))(tp.randn([4,3]), tp.randn([3])).shape)  # [4,3]

# randomness modes
try: tp.func.vmap(lambda: tp.randn([2]))(tp.randn([2]))
except RuntimeError as e: print("error mode:", "random" in str(e).lower())

print(tp.func.vmap(lambda: tp.randn([2]), randomness="same")(tp.randn([2])).shape)      # [2,2] with sharing
print(tp.func.vmap(lambda: tp.randn([2]), randomness="different")(tp.randn([2])).shape)  # [2,2] with independent draws

# vmap outranks autograd: grad inside vmap
def h(x): return (x * x).sum()
per_sample_grad = tp.func.vmap(tp.func.grad(h))(tp.randn([4,3]))
print(per_sample_grad.shape)  # [4,3]

# nested vmap
def dot(x, y): return (x * y).sum()
print(tp.func.vmap(tp.func.vmap(dot))(tp.randn([2,4,3]), tp.randn([2,4,3])).shape)  # [2,4]
PY
\end{lstlisting}

Expected shapes confirm that batch dims are placed at \texttt{out\_dims}, randomness modes branch correctly, and \texttt{vmap(grad(f))} produces per-sample gradients with an explicit graph (no silent fallback). A failure mode is also diagnostic: calling an unsupported op inside vmap throws \path{No native batch rule registered for <op>} from \path{gen\_api.py:624}, pointing directly at the missing entry in \path{BatchingKernels.cpp:1135}.
